\documentclass[../../../../hsm_tok_q.tex]{subfiles}

\begin{document}
    次の\:\myboxa{　　　}\:の中の
    「{\textsf あ}」「{\textsf い}」に当てはまる数字をそれぞれ答えよ。

    \textsf{図\ref{fig/hsm_tok_2010_2nd_01_08_01_q}}で，点Oは，線分ABを直径とする%
    円の中心であり，3点C，D，Eは円Oの周上にある点である．

    5点A，B，C，D，Eは，\textsf{図\ref{fig/hsm_tok_2010_2nd_01_08_01_q}}のように，%
    A，C，D，B，Eの順に並んでおり，互いに一致しない．

    点Aと点E，点Bと点C，点Bと点D，点Dと点Eをそれぞれ結び，%
    線分BCと線分DEとの交点をFとする．

    $\myarca{\mathrm{CD}}=\myfraca{1}{3}\myarca{\mathrm{AB}}$，%
    $\angle\mathrm{BAE}$＝$52^\circ$のとき，\\
    $x$で示した$\angle\mathrm{BFD}$の大きさは%
    \:\myboxa{\textsf{\hspace{.45\zw}\textsf{あ\hspace{.1\zw}い}\hspace{.45\zw}}}\:%
    度である．
    \\    
    \begin{minipage}{\linewidth}
        \vspace{.4\baselineskip}
        {\centering
            \setlength{\abovecaptionskip}{5pt}
            \includegraphics%
                {\subfix{fig_hsm_tok_2010_2nd_01_08/fig_hsm_tok_2010_2nd_01_08_01_q.pdf}}
                \captionof{figure}{} \label{fig/hsm_tok_2010_2nd_01_08_01_q}
        }
    \end{minipage}
    
    \vspace{.2\baselineskip}
\end{document}


